\section{A reversal reveals a missing constraint}\label{movinggoal}
The isolated theorem does not cover a request that changes while its old motion
is still active. A separate, predeclared comparison tests that case with the
same scalar plant and retained physical state. It uses three methods, four
request/deadline variants, three response rates and two constant disturbances:
72 parameter cases, each sampled at 1 ms and .5 ms. No parameters were tuned
after the results.

The first $20$ mm request is issued at $.20$ s and available at $.22$ s.
A reversal is issued at $.50$ s. All methods see identical requests and clocks.
The buffered case retains a common 60 ms execution commitment. Native replanning
uses its previous native state at execution time. FIR increments use the last
committed target. Neither process resets physical state.

\begin{center}\small
\begin{tabular}{@{}lrrr@{}}\toprule
Numerically valid tasks, out of six & Ruckig fast & Ruckig $.8$ s & FIR\\\midrule
Return to $0$; deadline $1.70$ s & 6 & 6 & 0\\
Same path; deadline before delivery & 0 & 0 & 0\\
Buffered return; deadline $1.90$ s & 6 & 6 & 6\\
Reverse to $-20$ mm; deadline $1.90$ s & 6 & 6 & 0\\\bottomrule
\end{tabular}
\end{center}

The task score requires at least 1 mm of initial progress, accepted requests,
causal service timing, nominal reference velocity/acceleration/jerk limits, and
the physical position, speed and acceleration bands from deadline through
$2.5$ s. Thus an always-hold policy cannot pass by ignoring the first move.
The common numerical reference-limit allowance is $10^{-10}$ in each respective
SI unit. This allowance is separate from native reconstruction tolerances;
small exact exported excesses remain recorded.

\textbf{The important failure.} In the unbuffered return, the exact FIR jerk peak
is $6.25\ \mathrm{m/s^3}$, below the $7.5$ limit. Its exact acceleration peak is
$21/32=.65625\ \mathrm{m/s^2}$, above the shared $.5$ limit. The jerk-only guard
correctly accepts its stated property, but that property is insufficient for
the full motion contract. Meeting the eventual endpoint does not cure this
limit violation. The buffered FIR return meets this numerical task contract.
The farther reversal fails its jerk admission check and remains a failed task.

The early-deadline case has the same reference and physical paths as the first
case. It changes only scoring. A deadline before command delivery violates this
service contract; it does not by itself prove universal physical impossibility.
All decisions agree at both sample steps. These are numerical task results,
not all-time task proofs. The earlier rounded-reference theorem is instantiated
for the twelve frozen gate06 paths, not these new native histories.

\textbf{Consequence.} The next useful admission result must include acceleration
and, where required, velocity and task time. A jerk-only proof is valuable but
cannot stand in for those conditions. This comparison provides no evidence of
general FIR superiority. Exact guard evidence, native fidelity checks, physical
traces are in \path{moving-goal-08/}; the independent review is in
\path{task-contract-08/grader/}.
